\defmodule {REDBAS2}

This module offers tools to reduce a basis in the
sense of Minkowski and find a shortest vector in a lattice.
The difference between this and {\tt REDBAS} is that here,
bounds on the numerical errors are computed to make sure that the
results are formally correct.

\bigskip\hrule
\code

DEFINITION MODULE REDBAS2;

FROM CONFIG   IMPORT BDim;
FROM LATBASIS IMPORT Basis;

PROCEDURE InitREDBAS2 (d : BDim);
\endcode
   \tab  Initializes the module. The value of {\tt d} should be at least
    as large as the dimension of the largest basis to be considered
    by the procedures below. \endtab
\code

PROCEDURE VerifMinkowski (B : Basis; d : BDim) : BOOLEAN;
\endcode
   \tab  Performs a formal verification and 
   returns {\tt TRUE} if the basis {\tt B} is really reduced. 
   Assumes that the first {\tt d} vectors are really reduced and
   verifies only the remaining ones. \endtab
\code

PROCEDURE ReductVerifMinkowski (B : Basis) : BOOLEAN;
\endcode
   \tab  Performs a reduction as in {\tt ReductMinkowski}, then verifies it
   formally, and if the basis is not really reduced, completes the reduction.
   This is usually much more costly than {\tt ReductMinkowski}, but the
   result is perfectly reliable. \endtab
\code

PROCEDURE VerifShortest (B : Basis) : BOOLEAN;
\endcode
   \tab  Performs a formal verification and 
   returns {\tt TRUE} if the first vector of basis {\tt B} 
   is the shortest vector in the lattice generated by {\tt B}.
   \endtab
\code

PROCEDURE ShortestVectorVerif (B : Basis) : BOOLEAN;
\endcode
   \tab  Computes the shortest vector as in {\tt ShortestVector},
   performs a formal verification, and if the vector was not really the 
   shortest due to numerical (floating-point) imprecisions, 
   computes the shortest one.
   Returns {\tt TRUE} if the procedure succeeds.
   This is usually much more costly than {\tt ShortestVector}, but the
   result is perfectly reliable.  \endtab
\code

END REDBAS2.
\endcode
